% ICLR 2027 migration of the final arXiv manuscript.
% Preliminaries and Experimental Setup are in Appendices A and B, respectively.
% Build from this directory: latexmk -pdf -interaction=nonstopmode -halt-on-error main.tex
\documentclass{article}
\PassOptionsToPackage{table}{xcolor}
\usepackage{iclr2027_conference,times}
\input{preamble}
\title{Kimodo-Policy: From Text-to-Motion Generators to Humanoid Vision-Language-Action Policies}
\input{author_metadata}

% Keep the official anonymous review mode for this working draft.
%\iclrfinalcopy

\begin{document}
Vision--language--action (VLA) models integrate visual perception, language understanding, and action generation across manipulation tasks and robot embodiments~\cite{rt2,openvla,pi0,wang2026limitsvisionandlanguagenavigation,wang2025dreamnavtrajectorybasedimaginativeframework}.
Learning instruction-conditioned behaviors from demonstrations provides a unified mapping from multimodal inputs to robot actions.
Extending VLA learning to humanoids could broaden task coverage in human environments, where locomotion, posture, and manipulation must be coordinated~\cite{grootn1,psi0,openhlm}.
For example, transferring objects between tables requires grasping, turning, walking, and placement while maintaining balance (Fig.~\ref{fig:overview}, bottom).
Such tasks require grounding task intent and visual feedback in coherent whole-body actions, motivating the extension of robot policy architectures to humanoid loco-manipulation.
\par\bigskip
Humanoid policy learning commonly extends the backbone--action expert paradigm from manipulation to whole-body control.
Pretrained vision--language models encode instruction and observation semantics (Fig.~\ref{fig:policy_routes}(a))~\cite{grootn1,psi0,openhlm}, while video generators supply spatiotemporal priors for scene evolution (Fig.~\ref{fig:policy_routes}(b))~\cite{motionwam,wholebodywam_wbc}.
These representations condition action generation across root, body, and hand motion.
Unified predictors generate whole-body sequences (Fig.~\ref{fig:policy_routes}(d))~\cite{openhlm,motionwam}; hybrid designs combine upper-body joint targets with lower-body locomotion commands executed by dedicated controllers (Fig.~\ref{fig:policy_routes}(e))~\cite{humanoidexo,wholebodyvla}.
The former captures cross-component dependencies; the latter couples upper-body control with learned locomotion and stabilization.
Both transfer pretrained semantic and visual dynamics knowledge to humanoid loco-manipulation.
\par\bigskip
Despite these advances, two challenges remain.
\textbf{\ding{182} Insufficient Whole-Body Motion Priors.}
Compared with fixed-base manipulation, humanoid control requires higher-dimensional inter-joint and temporal coordination. Backbone pretraining~\cite{grootn1,psi0,motionwam,wholebodywam_wbc} does not directly equip action generators with priors aligned with target humanoid kinematics, leaving them to learn these priors, often from limited task-specific demonstrations.
\textbf{\ding{183} Insufficient Differentiation of Motion Components.}
Root, body, and hand control differ in scale and precision. Unified prediction~\cite{openhlm,motionwam} may create competing objectives; hybrid execution~\cite{humanoidexo,wholebodyvla} does not itself provide dedicated root, body, and hand generators. Both must accommodate component-specific requirements while maintaining whole-body coordination.
\par\bigskip
\end{document}
